\section{Problem Formulation}

  FreeCell can be modeled as a deterministic, fully observable, single-agent search problem
  \[
  P = (\mathcal{S}, \mathcal{A}, T, s_0, G, c),
  \]
  where \(\mathcal{S}\) is the set of legal game states, \(\mathcal{A}\) is the set of legal actions, \(T\) is the transition model, \(s_0\) is the
  initial state, \(G\) is the goal test, and \(c\) is the step-cost function. Since all cards are visible from the beginning and no randomness appears
  after the initial deal, FreeCell is well suited for classical state-space search.

  \subsection{State Representation}

  A state \(n \in \mathcal{S}\) is represented as
  \[
  n = \bigl(\langle C_1, C_2, \dots, C_8 \rangle,\; \langle E_1, E_2, E_3, E_4 \rangle,\; \langle F_C, F_D, F_H, F_S \rangle \bigr),
  \]
  where:
  \begin{itemize}
      \item \(C_i\) is the \(i\)-th cascade (tableau column), represented as an ordered list of cards;
      \item \(E_j\) is the \(j\)-th free cell, which is either empty or contains exactly one card;
      \item \(F_s\) denotes the current top rank stored in the foundation of suit \(s \in \{C,D,H,S\}\), with \(0\) meaning that the foundation is
  empty.
  \end{itemize}

  Each card is represented by its \emph{rank} and \emph{suit}, i.e.,
  \[
  \text{card} = (r, s),
  \]
  where \(r \in \{1,\dots,13\}\) and \(s \in \{C,D,H,S\}\). In particular, \(r=1\) denotes Ace, \(r=11\) denotes Jack, \(r=12\) denotes Queen, and
  \(r=13\) denotes King.

  \subsection{Initial State}

  The initial state \(s_0\) is generated from a standard 52-card deck dealt face-up into 8 cascades. Specifically:
  \begin{itemize}
      \item 4 cascades contain 7 cards each;
      \item 4 cascades contain 6 cards each;
      \item all 4 free cells are empty;
      \item all 4 foundation piles are empty.
  \end{itemize}

  Formally, in the initial state,
  \[
  E_1 = E_2 = E_3 = E_4 = \emptyset
  \quad \text{and} \quad
  F_C = F_D = F_H = F_S = 0.
  \]

  \subsection{Actions (Transition Model)}

  From a state \(n\), a legal action corresponds to one valid move under the standard FreeCell rules. The transition model \(T(n,a)\) deterministically
  returns the unique successor state obtained after applying action \(a\).

  \begin{description}
      \item[Single-card moves.]
      A single top card from a cascade, or a card currently stored in a free cell, may be moved to:
      \begin{itemize}
          \item an empty free cell;
          \item another cascade, if it forms a descending alternating-color sequence with the destination top card;
          \item an empty cascade;
          \item its corresponding foundation, if it has the same suit and its rank is exactly one greater than the current top rank of that foundation.
      \end{itemize}

      \item[Sequence moves between cascades.]
      A sequence of cards may be moved from one cascade to another if:
      \begin{itemize}
          \item the moved sequence itself is a valid descending alternating-color sequence; 
          \item its length does not exceed the move capacity allowed by the current number of empty free cells and empty cascades.
      \end{itemize}
  \end{description}

  Let \(m\) be the number of empty free cells and let \(e\) be the number of empty cascades available as intermediate storage (excluding the
  destination cascade when it is empty). Then the maximum number of cards that can be moved together is
  \[
  ( m + 1 ) 2^e.
  \]

  Thus, a sequence move of length \(k\) is legal only if
  \[
  k \le (m+1)2^e.
  \]

  \subsection{Goal Test}

  A state \(n\) is a goal state if all cards have been moved to the foundations. Equivalently,
  \[
  G(n) = \text{true}
  \quad \Longleftrightarrow \quad
  F_C = F_D = F_H = F_S = 13.
  \]
  This is equivalent to saying that the total number of cards in the foundations is 52:
  \[
  F_C + F_D + F_H + F_S = 52.
  \]

  \subsection{Path Cost}

\subsubsection{Standard Unit-Cost Model}
In the general formulation of FreeCell, we primarily adopt a \textbf{unit-cost model}. Every legal move—including single-card moves and valid sequence moves—is assigned a cost of 1. Thus, if a path from the initial state $s_0$ to a state $n$ consists of $d$ legal moves, the path cost is defined as:
\[
g(n) = d.
\]
Equivalently, for any action $a$:
\[
c(n,a,n') = 1.
\]
Under this formulation, $g(n)$ measures the total number of moves made so far. To remain consistent with this model, any heuristic $h(n)$ used in informed search (such as $A^*$) estimates the remaining cost in the same unit: the number of moves. This model ensures that BFS, IDS, and $A^*$ find the \textbf{optimal solution} in terms of the minimum number of steps.

\subsubsection{Weighted Strategic Cost Model (UCS Implementation)}
To guide the \textbf{Uniform-Cost Search (UCS)} more efficiently through the state space, we implement a \textbf{weighted cost model} as detailed in Algorithm 7 and the source code \texttt{ucs.py}. This model assigns dynamic costs based on the strategic value of each move:

\begin{table}[h]
\centering
\begin{tabular}{|l|c|l|}
\hline
\textbf{Move Type} & \textbf{Cost} & \textbf{Strategic Reasoning} \\ \hline
To Foundation & 1 & Highest priority; directly advances toward the goal. \\ \hline
Emptying a Cascade & 5 & High priority; creates empty columns for better mobility. \\ \hline
Cascade to Cascade & 10 & Standard move; default cost for tableau maneuvers. \\ \hline
To Free Cell & 50 & Penalty cost; discourages filling limited storage slots. \\ \hline
\end{tabular}
\caption{Dynamic move costs used in UCS expansion.}
\end{table}

While this weighted model might not minimize the absolute number of moves, it significantly reduces the number of expanded nodes by prioritizing strategically "productive" paths over stalling maneuvers.